% GENERATED FILE. Edit canonical lessons, then run npm run generate:books.
% canonical-source-hash: fnv1a64:b1289cc5797ec994
% canonical-lessons: HI-C480-svad, HI-C480-spars, HI-C480-gandh, HI-C480-mahak, HI-C480-catna

\chapter{The Senses}
\label{ch:hi-the-senses}
\hlchaptermodality{480}

\begin{chapteropening}
\textbf{By the end of this chapter:} \emph{I can say how something tastes, smells or feels.}

Complete the last lesson of chapter 480: I can say how something tastes, smells or feels.
\end{chapteropening}

\section[\texorpdfstring{\hi{स्वाद} (svād)}{स्वाद (svād)}]{\hi{स्वाद} (svād) — a taste; a flavour}

\label{lesson:HI-C480-svad}



\begin{quote}
Before the new one: say the Hindi for to chew, then the Hindi for to spit.
\end{quote}

\subsection*{You'll want to know: \hi{स्वाद}}
\textbf{\hi{स्वाद}} — \emph{svād} — \textquotedblleft{}a taste; a flavour\textquotedblright{}.

\begin{quote}
\hi{माँ} \hi{के} \hi{हाथ} \hi{के} \hi{खाने} \hi{का} \hi{स्वाद} \hi{सबसे} \hi{अच्छा} \hi{है।}
\end{quote}

\begin{quote}
\emph{māṁ ke hāth ke khāne kā svād sabse acchā hai.}
\end{quote}

\begin{quote}
Food cooked by my mother tastes best of all.
\end{quote}

Masculine.

\begin{cousinweb}
\textbf{\hi{स्वाद}} (\emph{svād}) is a tatsama, borrowed from Sanskrit \emph{svāda}, \textquotedblleft{}taste, flavour, savour\textquotedblright{} (Monier-Williams).
\end{cousinweb}

\begin{grammarlens}[title={What you've built}]
One more word for describing what you lived and felt.
\end{grammarlens}

\subsection*{Your turn}
\begin{itemize}
\raggedright
  \item \emph{Say it:} \emph{svād}
  \item \emph{Say it:} \emph{svād}, once more
  \item \emph{Say it:} \emph{māṁ ke hāth ke khāne kā svād sabse acchā hai.} — Food cooked by my mother tastes best of all.
  \item \emph{Recall:} say the Hindi for to chew, then the Hindi for to spit, then say \emph{svād} again
\end{itemize}

\begin{tcolorbox}[breakable,colback=teal!4,colframe=teal!35!black,title={Before you move on}]
What does \hi{स्वाद} mean? (\textquotedblleft{}A taste; a flavour\textquotedblright{}.) Say it once more.
\end{tcolorbox}
\section[\texorpdfstring{\hi{स्पर्श} (sparś)}{स्पर्श (sparś)}]{\hi{स्पर्श} (sparś) — a touch}

\label{lesson:HI-C480-spars}



\begin{quote}
Before the new one: say the Hindi for to spit, then the Hindi for a taste.
\end{quote}

\subsection*{You'll want to know: \hi{स्पर्श}}
\textbf{\hi{स्पर्श}} — \emph{sparś} — \textquotedblleft{}a touch\textquotedblright{}.

\begin{quote}
\hi{माँ} \hi{के} \hi{स्पर्श} \hi{से} \hi{बच्चा} \hi{सो} \hi{गया।}
\end{quote}

\begin{quote}
\emph{māṁ ke sparś se baccā so gayā.}
\end{quote}

\begin{quote}
At its mother's touch the baby fell asleep.
\end{quote}

Masculine.

\begin{cousinweb}
\textbf{\hi{स्पर्श}} (\emph{sparś}) is a tatsama, borrowed from Sanskrit \emph{sparśa}, \textquotedblleft{}touch, the sense of touch\textquotedblright{} (Monier-Williams). It is the noun; the everyday verb is \emph{chūnā}.
\end{cousinweb}

\begin{grammarlens}[title={What you've built}]
One more word for describing what you lived and felt.
\end{grammarlens}

\subsection*{Your turn}
\begin{itemize}
\raggedright
  \item \emph{Say it:} \emph{sparś}
  \item \emph{Say it:} \emph{sparś}, once more
  \item \emph{Say it:} \emph{māṁ ke sparś se baccā so gayā.} — At its mother's touch the baby fell asleep.
  \item \emph{Recall:} say the Hindi for to spit, then the Hindi for a taste, then say \emph{sparś} again
\end{itemize}

\begin{tcolorbox}[breakable,colback=teal!4,colframe=teal!35!black,title={Before you move on}]
What does \hi{स्पर्श} mean? (\textquotedblleft{}A touch\textquotedblright{}.) Say it once more.
\end{tcolorbox}
\section[\texorpdfstring{\hi{गंध} (gandh)}{गंध (gandh)}]{\hi{गंध} (gandh) — a smell; an odour}

\label{lesson:HI-C480-gandh}



\begin{quote}
Before the new one: say the Hindi for a taste, then the Hindi for a touch.
\end{quote}

\subsection*{You'll want to know: \hi{गंध}}
\textbf{\hi{गंध}} — \emph{gandh} — \textquotedblleft{}a smell; an odour\textquotedblright{}.

\begin{quote}
\hi{रसोई} \hi{से} \hi{मसालों} \hi{की} \hi{गंध} \hi{आ} \hi{रही} \hi{है।}
\end{quote}

\begin{quote}
\emph{rasoī se masāloṁ kī gandh ā rahī hai.}
\end{quote}

\begin{quote}
The smell of spices is coming from the kitchen.
\end{quote}

Feminine.

\begin{cousinweb}
\textbf{\hi{गंध}} (\emph{gandh}) is inherited: CDIAL lists Hindi \emph{gandh}, \textquotedblleft{}smell\textquotedblright{}, under Sanskrit \emph{gandhá}, \textquotedblleft{}smell\textquotedblright{} (CDIAL 4014). It can be pleasant or not; \emph{khuśbū} is only a good smell.
\end{cousinweb}

\begin{grammarlens}[title={What you've built}]
One more word for describing what you lived and felt.
\end{grammarlens}

\subsection*{Your turn}
\begin{itemize}
\raggedright
  \item \emph{Say it:} \emph{gandh}
  \item \emph{Say it:} \emph{gandh}, once more
  \item \emph{Say it:} \emph{rasoī se masāloṁ kī gandh ā rahī hai.} — The smell of spices is coming from the kitchen.
  \item \emph{Recall:} say the Hindi for a taste, then the Hindi for a touch, then say \emph{gandh} again
\end{itemize}

\begin{tcolorbox}[breakable,colback=teal!4,colframe=teal!35!black,title={Before you move on}]
What does \hi{गंध} mean? (\textquotedblleft{}A smell; an odour\textquotedblright{}.) Say it once more.
\end{tcolorbox}
\section[\texorpdfstring{\hi{महक} (mahak)}{महक (mahak)}]{\hi{महक} (mahak) — an aroma; a scent}

\label{lesson:HI-C480-mahak}



\begin{quote}
Before the new one: say the Hindi for a touch, then the Hindi for a smell.
\end{quote}

\subsection*{You'll want to know: \hi{महक}}
\textbf{\hi{महक}} — \emph{mahak} — \textquotedblleft{}an aroma; a scent\textquotedblright{}.

\begin{quote}
\hi{बारिश} \hi{के} \hi{बाद} \hi{मिट्टी} \hi{की} \hi{महक} \hi{आती} \hi{है।}
\end{quote}

\begin{quote}
\emph{bāriś ke bād miṭṭī kī mahak ātī hai.}
\end{quote}

\begin{quote}
After the rain there is a scent of earth.
\end{quote}

Feminine.

\begin{cousinweb}
\textbf{\hi{महक}} (\emph{mahak}) is a tadbhava: CDIAL lists Hindi \emph{mahak}, \textquotedblleft{}fragrance\textquotedblright{}, under Sanskrit \emph{magha}, which it glosses as a fragrance and the name of a flower (CDIAL 9703). It is a strong, spreading smell, usually a pleasant one.
\end{cousinweb}

\begin{grammarlens}[title={What you've built}]
One more word for describing what you lived and felt.
\end{grammarlens}

\subsection*{Your turn}
\begin{itemize}
\raggedright
  \item \emph{Say it:} \emph{mahak}
  \item \emph{Say it:} \emph{mahak}, once more
  \item \emph{Say it:} \emph{bāriś ke bād miṭṭī kī mahak ātī hai.} — After the rain there is a scent of earth.
  \item \emph{Recall:} say the Hindi for a touch, then the Hindi for a smell, then say \emph{mahak} again
\end{itemize}

\begin{tcolorbox}[breakable,colback=teal!4,colframe=teal!35!black,title={Before you move on}]
What does \hi{महक} mean? (\textquotedblleft{}An aroma; a scent\textquotedblright{}.) Say it once more.
\end{tcolorbox}
\section[\texorpdfstring{\hi{चाटना} (cāṭnā)}{चाटना (cāṭnā)}]{\hi{चाटना} (cāṭnā) — to lick}

\label{lesson:HI-C480-catna}



\begin{quote}
Before the new one: say the Hindi for a smell, then the Hindi for an aroma.
\end{quote}

\subsection*{You'll want to know: \hi{चाटना}}
\textbf{\hi{चाटना}} — \emph{cāṭnā} — \textquotedblleft{}to lick\textquotedblright{}.

\begin{quote}
\hi{बच्चे} \hi{ने} \hi{उँगली} \hi{से} \hi{शहद} \hi{चाटा।}
\end{quote}

\begin{quote}
\emph{bacce ne uṅglī se śahad cāṭā.}
\end{quote}

\begin{quote}
The child licked honey off his finger.
\end{quote}

Past, with \hi{ने} (he or she …): \hi{उसने} \hi{शहद} \hi{चाटा}, \emph{usne śahad cāṭā}.

\begin{cousinweb}
\textbf{\hi{चाटना}} (\emph{cāṭnā}) is a tadbhava: CDIAL lists Hindi \emph{cāṭnā}, \textquotedblleft{}to lick, taste, eat\textquotedblright{}, under \emph{*caṭṭ}, a form it reconstructs (CDIAL 4573).
\end{cousinweb}

\begin{grammarlens}[title={What you've built}]
One more word for describing what you lived and felt.
\end{grammarlens}

\subsection*{Your turn}
\begin{itemize}
\raggedright
  \item \emph{Say it:} \emph{cāṭnā}
  \item \emph{Say it:} \emph{cāṭnā}, once more
  \item \emph{Say it:} \emph{bacce ne uṅglī se śahad cāṭā.} — The child licked honey off his finger.
  \item \emph{Recall:} say the Hindi for a smell, then the Hindi for an aroma, then say \emph{cāṭnā} again
\end{itemize}

\begin{tcolorbox}[breakable,colback=teal!4,colframe=teal!35!black,title={Before you move on}]
What does \hi{चाटना} mean? (\textquotedblleft{}To lick\textquotedblright{}.) Say it once more.
\end{tcolorbox}
